\subsection{Purpose}

The purpose of this section is to establish minimum requirements for preparing for, identifying, reporting, assessing, containing, investigating, recovering from, and learning from security and technology incidents.

These requirements are intended to:
\begin{itemize}
    \item Reduce the impact and duration of security and operational incidents
    \item Ensure that suspected incidents are reported and handled consistently
    \item Establish clear responsibilities and communication requirements
    \item Protect systems, data, users, and organizational operations
    \item Preserve relevant evidence for investigation and follow-up
    \item Support timely recovery and restoration of affected services
    \item Identify corrective actions that reduce the likelihood of recurrence
\end{itemize}

\subsection{Incident Definition}

An incident is an actual or suspected event that may negatively affect the confidentiality, integrity, or availability of systems, services, accounts, networks, or data.

Incidents may include:
\begin{itemize}
    \item Unauthorized access or attempted unauthorized access
    \item Malware, ransomware, or suspected system compromise
    \item Phishing, credential theft, or account takeover
    \item Loss or unauthorized disclosure of sensitive information
    \item Lost or stolen devices
    \item Accidental deletion, alteration, or destruction of data
    \item Denial-of-service or significant service interruption
    \item Unauthorized changes to systems, applications, or network controls
    \item Backup failure affecting recovery capability
    \item Physical damage or environmental conditions affecting critical systems
    \item Violations of security procedures or approved configurations
\end{itemize}

Personnel \must{} report suspected incidents even when the facts are incomplete or the impact is not yet known.

\subsection{Incident Response Responsibilities}

Management \must{} designate a person or team responsible for coordinating incident response. The incident-response function \should{} include, where appropriate:
\begin{itemize}
    \item Incident coordinator
    \item System and network administrators
    \item System and data owners
    \item Management representatives
    \item Legal, compliance, or privacy representatives
    \item Communications or public-relations representatives
    \item Human resources representatives
    \item Relevant vendors, insurers, or external specialists
\end{itemize}

Incident-response responsibilities \must{} be assigned before an incident occurs. Contact information \must{} be maintained and reviewed periodically.

Personnel involved in incident response \should{} receive appropriate training and participate in exercises or testing.

\subsection{Incident Reporting}

Users and administrators \must{} report suspected incidents promptly to management, the designated security contact, or another approved reporting channel.

Incident reports \should{} include:
\begin{itemize}
    \item Date and time of discovery
    \item Person reporting the event
    \item Affected user, device, account, system, service, or data
    \item Description of the observed activity
    \item Actions already taken
    \item Available logs, messages, screenshots, or other evidence
    \item Known or suspected impact
\end{itemize}

Personnel \must{} not delay reporting while attempting to determine the complete cause or scope of an event.

Incident reports and response communications \must{} be handled according to their sensitivity. Information about an active incident \must{} be shared only with personnel who have a business or response-related need to know.

\subsection{Incident Classification}

Incidents \must{} be classified according to their severity, scope, affected resources, data involved, operational impact, and potential legal or contractual consequences.

Classification \should{} consider:
\begin{itemize}
    \item Number and importance of affected systems
    \item Type and sensitivity of affected information
    \item Whether unauthorized access is confirmed or suspected
    \item Availability and operational impact
    \item Duration and current status of the event
    \item Likelihood of continued or repeated activity
    \item Potential regulatory, contractual, legal, or safety impact
\end{itemize}

The incident coordinator \must{} assign or confirm the severity level and \must{} escalate the incident when new information indicates greater impact or risk.

\subsection{Incident Response Process}

Incident response \must{} follow a documented process that includes:
\begin{enumerate}
    \item Preparation
    \item Identification and reporting
    \item Initial assessment and classification
    \item Containment
    \item Evidence preservation and investigation
    \item Eradication
    \item Recovery and service restoration
    \item Post-incident review and corrective action
\end{enumerate}

The order and extent of response activities may vary according to the incident. Actions \must{} be coordinated to avoid increasing the damage, destroying evidence, or disrupting critical operations unnecessarily.

\subsection{Containment}

The incident-response team \must{} take reasonable steps to contain an incident and prevent additional harm.

Containment actions may include:
\begin{itemize}
    \item Isolating affected devices or network segments
    \item Disabling compromised accounts
    \item Revoking sessions, tokens, keys, or credentials
    \item Blocking malicious addresses, domains, files, or processes
    \item Removing unauthorized access
    \item Suspending affected services
    \item Restricting access to affected systems or data
    \item Activating alternate systems or recovery procedures
\end{itemize}

Containment actions \must{} consider business continuity, safety, evidence preservation, and the possibility that the attacker or malware remains active.

Compromised credentials \must{} be changed or revoked. Replacement credentials \must{} be generated and distributed securely.

\subsection{Evidence Preservation and Investigation}

Relevant evidence \must{} be preserved when an incident may require investigation, legal action, regulatory reporting, insurance notification, or other formal review.

Evidence may include:
\begin{itemize}
    \item System, application, network, authentication, and security logs
    \item Disk, memory, or system images where appropriate
    \item Malware samples or suspicious files
    \item Email messages, headers, links, and attachments
    \item Screenshots and error messages
    \item Access records and account activity
    \item Configuration files and change records
    \item Relevant physical security or access records
\end{itemize}

Evidence \must{} be protected from alteration or unauthorized access. Investigators \should{} document when evidence was collected, by whom, from which source, and how it was stored or transferred.

Administrators \mustnot{} alter or delete relevant logs, files, accounts, or systems without considering evidence-preservation requirements and coordinating with the incident coordinator.

\subsection{Eradication and Recovery}

After containment and investigation, the response team \must{} remove the cause of the incident and address known weaknesses before returning affected systems to normal operation.

Eradication and recovery activities may include:
\begin{itemize}
    \item Removing malware or unauthorized software
    \item Rebuilding or restoring compromised systems
    \item Applying security patches and configuration changes
    \item Resetting affected passwords and authentication factors
    \item Removing unauthorized accounts, rules, services, or access paths
    \item Restoring data from verified backups
    \item Testing restored systems before returning them to production
    \item Increasing monitoring for signs of continued compromise
\end{itemize}

Systems \must{} be validated before being returned to normal use. Validation \should{} confirm that the system is secure, functional, appropriately monitored, and connected only to authorized services.

Recovery activities \must{} be coordinated with system owners and business stakeholders. Critical services \should{} be restored according to documented recovery priorities.

\subsection{Communication and Notification}

Incident communications \must{} be coordinated by authorized personnel. Personnel \mustnot{} make public statements, contact affected external parties, or notify authorities unless authorized to do so.

Management \must{} determine whether an incident requires notification to:
\begin{itemize}
    \item Affected individuals
    \item Customers or business partners
    \item Vendors or service providers
    \item Insurers
    \item Regulators or government authorities
    \item Law enforcement
\end{itemize}

Required notifications \must{} be made within applicable time limits and using approved communication methods. Incident communications \should{} be accurate, limited to confirmed information, and updated when significant facts change.

\subsection{Post-Incident Review}

A post-incident review \must{} be performed for significant incidents and \should{} be performed for recurring or high-risk events.

The review \should{} identify:
\begin{itemize}
    \item What happened and when
    \item Systems, accounts, data, and services affected
    \item How the incident was detected
    \item Actions that were successful or unsuccessful
    \item Root or contributing causes
    \item Control failures or process weaknesses
    \item Required corrective actions
    \item Responsible owners and target completion dates
    \item Changes needed to policies, procedures, training, or technology
\end{itemize}

Corrective actions \must{} be tracked until completed or formally accepted as a risk. Lessons learned \should{} be incorporated into security controls, incident procedures, training, and recovery plans.

\subsection{Incident Records}

An incident record \must{} be maintained for each reported incident. The record \should{} include:
\begin{itemize}
    \item Incident identifier
    \item Date and time reported
    \item Reporter and incident coordinator
    \item Classification and severity
    \item Affected systems, accounts, services, and data
    \item Timeline of significant events and response actions
    \item Evidence collected and its storage location
    \item Notifications and communications
    \item Recovery actions
    \item Post-incident findings
    \item Corrective actions and completion status
    \item Date the incident was closed
\end{itemize}

Incident records \must{} be protected from unauthorized access and retained according to applicable legal, contractual, business, and records-management requirements.

\subsection{Testing and Maintenance}

The incident-response process \must{} be reviewed periodically and after significant incidents, organizational changes, or changes to critical systems.

Incident-response procedures \should{} be tested through exercises, walkthroughs, simulations, or technical testing. Testing \should{} include relevant personnel, systems, communication methods, escalation paths, backups, and recovery procedures.

Identified weaknesses \must{} be documented and assigned for correction. Contact lists, response procedures, system inventories, recovery information, and escalation instructions \must{} be kept current.